%%%%%%%%%%%%%%%%%%%%%%%%%%%%%%%%%
%
%       EXERCÍCIO
%
%%%%%%%%%%%%%%%%%%%%%%%%%%%%%%%%%

\ifdefstring{\atividade}{prova}{%
    \ifdefstring{\modo}{objetivo}{%
        \renewcommand{\valorquestao}{\ValorQObj\ ponto}
    }{%
        \renewcommand{\valorquestao}{\ValorQDisc\ pontos}
    }%
}%

\begin{exercicioBanco}[\valorquestao]
Determine quais das funções a seguir são funções afins.
%
\begin{center}
    \begin{tabularx}{\textwidth}{XXXX}
        (I) \(f(x) = 0\). & (II) \(f(x) = x^{2}\). & (III) \(f(x) = \pi x - \sqrt{2}\). & (IV) \(f(x) = 2^{x}\).  \\[10pt]
        (V) \(f(x) = x\). & (VI) \(f(x) = -\dfrac{1}{2} - 2x\). & (VII) \(f(x) = \dfrac{1}{x}\). & (VIII) \(f(x) = \sqrt{x}\).
    \end{tabularx}
\end{center}

% Define as alternativas
\newcommand{\alternativas}{%
    \begin{center}
        \begin{tabularx}{\textwidth}{XXXXX}
            (a) I, III, V, VI. &
            (b) V, VI. &
            (c) V, VI, VII. &
            (d) Todas. &
            (e) Todas, exceto IV.
        \end{tabularx}
    \end{center}
}

% Define a resposta correta
\newcommand{\resposta}{A}

% Lógica condicional para exibição
\ifdefstring{\atividade}{lista}{%
    \alternativas
    \vspace{0.5em}
    
    \noindent\textbf{Resposta:} letra \textbf{\resposta}.
}{%
    \ifdefstring{\modo}{objetivo}{%
        \alternativas
    }{%
        % modo = discursiva → não mostra alternativas
    }
}
\end{exercicioBanco}

